\def\handoutsay{false}\def\handoutfill{false}
\documentclass[12pt]{article}
\newcommand{\handoutrunning}{Lecture 9: Covariance matrices and principal components}
\providecommand{\handoutsay}{true}
\usepackage[T1]{fontenc}
\usepackage{lmodern,microtype}
\usepackage[a4paper,margin=15mm,headheight=14pt,headsep=5mm]{geometry}
\usepackage{amsmath,amssymb,amsthm,mathtools}
\usepackage{enumitem,booktabs,needspace}
\usepackage{pifont}
\usepackage[dvipsnames]{xcolor}
\usepackage{fancyhdr}
\usepackage{ifthen}
\usepackage[hidelinks]{hyperref}
\hypersetup{pdfauthor={Chang-Ye Tu}}
\pagestyle{fancy}
\fancyhf{}
\fancyhead[L]{\small Linear Algebra}
\fancyhead[R]{\small\handoutrunning}
\fancyfoot[C]{\small\thepage}
\renewcommand{\headrulewidth}{0.3pt}
\setlength{\parindent}{0pt}
\setlength{\parskip}{4pt plus 1pt minus 1pt}
\setlength{\emergencystretch}{1.5em}
\usepackage{etoolbox}
\newcommand{\handoutdisplayskips}{\setlength{\abovedisplayskip}{7pt plus 2pt minus 2pt}%
  \setlength{\belowdisplayskip}{7pt plus 2pt minus 2pt}}
\apptocmd{\normalsize}{\handoutdisplayskips}{}{\PackageWarning{handout}{could not patch normalsize}}
\handoutdisplayskips
\setlist[enumerate]{label=(\alph*),leftmargin=*,itemsep=2pt,topsep=3pt}
\setlist[itemize]{leftmargin=*,itemsep=2pt,topsep=3pt}
\newcommand{\R}{\mathbb R}
\newcommand{\C}{\mathbb C}
\newcommand{\iu}{\mathrm{i}}
\newcommand{\contradiction}{\mathord{\text{\large\ding{56}}}}
\newcommand{\one}{\mathbf 1}
\DeclareMathOperator{\tr}{tr}
\DeclareMathOperator{\col}{col}
\DeclareMathOperator{\diag}{diag}
\DeclareMathOperator{\rank}{rank}
\DeclareMathOperator\prb{\mathsf{P}}
\DeclareMathOperator\expc{\mathsf{E}}
\DeclareMathOperator\var{var}
\DeclareMathOperator\cov{cov}
\DeclareMathOperator\cor{cor}
\DeclareMathOperator\argmin{argmin}
\DeclareMathOperator\argmax{argmax}
\DeclareMathOperator{\Span}{span}
\theoremstyle{definition}
\newtheorem{theorem}{Theorem}
\newtheorem{proposition}[theorem]{Proposition}
\newtheorem{lemma}[theorem]{Lemma}
\newtheorem{corollary}[theorem]{Corollary}
\newtheorem{definition}[theorem]{Definition}
\newtheorem{example}[theorem]{Example}
\newtheorem{remark}[theorem]{Remark}
\newtheorem{fact}{Fact}
\renewcommand{\thefact}{F\arabic{fact}}
\AddToHook{cmd/subsection/before}{\par\Needspace{9\baselineskip}}
\AddToHook{env/theorem/before}{\par\Needspace{4\baselineskip}}
\AddToHook{env/proposition/before}{\par\Needspace{4\baselineskip}}
\AddToHook{env/lemma/before}{\par\Needspace{4\baselineskip}}
\AddToHook{env/definition/before}{\par\Needspace{4\baselineskip}}
\AddToHook{env/example/before}{\par\Needspace{4\baselineskip}}
\definecolor{saybg}{gray}{0.86}
\NewDocumentEnvironment{say}{+b}
  {\ifthenelse{\equal{\handoutsay}{true}}{\par\Needspace{2\baselineskip}\noindent\colorbox{saybg}{\parbox{\dimexpr\linewidth-2\fboxsep\relax}%
   {\small\sffamily\ding{43}\enspace\textbf{Say.}\enspace#1}}\par}{}}
  {}
\providecommand{\handoutfill}{false}
\newsavebox{\boardbox}
\newlength{\boardht}
\definecolor{boardrule}{gray}{0.55}
\makeatletter
\newcommand{\boardstart}{\par\if@inlabel\leavevmode\par\fi}
\makeatother
\NewDocumentEnvironment{board}{o +b}
  {\boardstart\IfValueT{#1}{\noindent#1\par\nopagebreak}%
   \noindent\begin{lrbox}{\boardbox}\begin{minipage}{\linewidth}#2\end{minipage}\end{lrbox}%
   \setlength{\boardht}{\dimexpr\ht\boardbox+\dp\boardbox\relax}%
   \ifthenelse{\equal{\handoutfill}{true}}%
     {\setlength{\boardht}{1.4\boardht}%
      {\setlength{\fboxsep}{0pt}\setlength{\fboxrule}{0.4pt}\color{boardrule}%
       \framebox[\dimexpr\linewidth-2\fboxrule\relax]{\rule{0pt}{\boardht}}}}%
     {\ifthenelse{\equal{\handoutsay}{true}}%
        {\llap{\color{boardrule}\rule[-\dp\boardbox]{1.5pt}{\boardht}\hspace{4pt}}}{}%
      \usebox{\boardbox}}%
   \par}
  {}
\newcommand{\blank}[1]{\ifthenelse{\equal{\handoutfill}{true}}{\underline{\phantom{#1}}}{#1}}
\newcommand{\lec}[1]{Lecture~#1}
\newcommand{\handouttitle}[2]{%
  \begin{center}
  {\LARGE\sffamily\bfseries #2\par}\vspace{3pt}
  {\sffamily Linear Algebra \textperiodcentered\ \lec{#1}\par}
  \ifthenelse{\equal{\handoutsay}{true}}{\vspace{3pt}{\small\sffamily\colorbox{saybg}{\ding{43}\enspace\textbf{Say.}}\enspace
    teacher's copy: the shaded boxes are not in the students' copy.\par}}{}
  \ifthenelse{\equal{\handoutfill}{true}}{\vspace{3pt}{\small\sffamily Fill-in copy: complete the framed spaces in
    class.\par}}{}
  \end{center}}

\hypersetup{pdftitle={Linear Algebra, Lecture 9: Covariance matrices and principal components}}
\newcommand{\boardtop}{\vspace*{-\abovedisplayskip}}
\DeclareMathOperator{\row}{row}
\begin{document}
\handouttitle{9}{Covariance matrices and principal components}

\textbf{Last time} (\lec{8}): symmetric \(A=Q\Lambda Q'\) (\(Q\) orthogonal, \(\Lambda\) diagonal); PSD \(\iff\)
eigenvalues \(\geqslant0\). Earlier results are cited where they are used; F1, F2: probability, without proof.
\(I=I_n\), \(e_j=\col_j(I)\).

\textbf{Facts} on random variables (numbers whose values depend on chance; here with finite variances).
\begin{enumerate}[label=\textbf{F\arabic*.},ref=F\arabic*,leftmargin=*,itemsep=0pt,parsep=0pt,topsep=2pt]
\item\label{f-prob} \(\var(u)=\expc\bigl((u-\expc(u))^2\bigr)\geqslant0\),
\(\cov(u,v)=\expc\bigl((u-\expc(u))(v-\expc(v))\bigr)=\cov(v,u)\), \(\var(u)=\cov(u,u)\).
\item\label{f-lin} \(\expc(au+bv+c)=a\expc(u)+b\expc(v)+c\) for numbers \(a\), \(b\), \(c\) (also for finite sums).
\end{enumerate}

\subsection*{1. Random vectors and covariance matrices}

\begin{definition}\label{h-def-random}
A \textbf{random vector} \(z=(z_1,\dots,z_n)'\) or \textbf{random matrix} \(Z=(z_{ij})\) has random variables as
entries; \(\expc(Z)=(\expc(z_{ij}))\). The \textbf{mean vector} of \(z\) is
\(\mu=\expc(z)\), and its \textbf{covariance matrix} is
\[
\Sigma=\var(z)=\expc\bigl((z-\mu)(z-\mu)'\bigr)\in\R^{n\times n}
\]
By \lec{1}, Theorem~4(d) and \ref{f-prob}, \(\sigma_{ij}=\expc\bigl((z_i-\mu_i)(z_j-\mu_j)\bigr)=\cov(z_i,z_j)\),
\(\sigma_{ii}=\var(z_i)\), and \(\Sigma\) is symmetric.
\end{definition}

\begin{definition}\label{h-def-cov}
For random vectors \(u\in\R^p\), \(v\in\R^q\) with \(\mu_u=\expc(u)\), \(\mu_v=\expc(v)\), the
\textbf{covariance matrix of \(u\) and \(v\)} is
\[
\cov(u,v)=\expc\bigl((u-\mu_u)(v-\mu_v)'\bigr)\in\R^{p\times q}
\]
Its \((i,j)\) entry is \(\cov(u_i,v_j)\). So \(\cov(u,u)=\var(u)\) and
\(\cov(v,u)=\cov(u,v)'\) (\lec{1}, Theorem~8). \(u\) and \(v\) are \textbf{uncorrelated} if
\(\cov(u,v)=0\).
\end{definition}

\Needspace{17\baselineskip}
\begin{theorem}\label{h-thm-linear}
Let \(z\in\R^n\) be a random vector, \(\expc(z)=\mu\), \(\var(z)=\Sigma\).
\begin{enumerate}
\item \(\expc(AZB+D)=A\expc(Z)B+D\) for a random \(Z\in\R^{m\times p}\) and constant
\(A\in\R^{q\times m}\), \(B\in\R^{p\times l}\), \(D\in\R^{q\times l}\).
\item \(\expc(Az+b)=A\mu+b\) and \(\var(Az+b)=A\Sigma A'\) for constant \(A\in\R^{m\times n}\),
\(b\in\R^m\).
\item \(a'\Sigma a=\var(a'z)\geqslant0\) for \(a\in\R^n\): \(\Sigma\) is PSD.
\end{enumerate}
\end{theorem}
\begin{proof}
(a) Entry \((i,j)\) of \(AZB+D\) is \(\sum_r\sum_sa_{ir}z_{rs}b_{sj}+d_{ij}\) (\lec{1}, Theorem~4(d)); its expectation is
\(\sum_r\sum_sa_{ir}\expc(z_{rs})b_{sj}+d_{ij}\) (\ref{f-lin}), entry \((i,j)\) of \(A\expc(Z)B+D\).
\begin{board}[(b):]
\boardtop
\begin{align*}
\expc(Az+b)&=A\mu+b,\qquad Az+b-\expc(Az+b)=A(z-\mu)&&\text{((a): \(Z=z\), \(B=1\), \(D=b\))}\\
\var(Az+b)&=\expc\bigl(A(z-\mu)(z-\mu)'A'\bigr)=A\Sigma A'&&\text{(\lec{1}, Theorem~8; (a))}
\end{align*}
\end{board}
\begin{board}[(c):]
\boardtop
\begin{align*}
\var(a'z)&=a'\Sigma a\geqslant0&&\text{((b) with \(a'\), \(b=0\); \ref{f-prob})}
\end{align*}
\(\Sigma\) is symmetric (Definition~\ref{h-def-random}): PSD (\lec{8}, Definition~13: \(x'\Sigma x\geqslant0\) for
all \(x\)).
\end{board}
\nopagebreak
\end{proof}

\Needspace{15\baselineskip}
\begin{theorem}\label{h-thm-cov}
Let \(u\in\R^p\), \(v\in\R^q\) be random vectors, \(\mu_u=\expc(u)\), \(\mu_v=\expc(v)\).
\begin{enumerate}
\item \(\cov(Au+b,Bv+d)=A\cov(u,v)B'\) for constant \(A\in\R^{k\times p}\), \(b\in\R^k\),
\(B\in\R^{l\times q}\), \(d\in\R^l\). So \(\cov(Az,Bz)=A\Sigma B'\) if \(\var(z)=\Sigma\).
\item \(\cov(u_1+u_2,v)=\cov(u_1,v)+\cov(u_2,v)\) and
\(\cov(v,u_1+u_2)=\cov(v,u_1)+\cov(v,u_2)\) (\(u_1,u_2\in\R^p\)).
\item If \(x_1,\dots,x_m\in\R^p\) and \(\cov(x_i,x_j)=0\) for \(i\ne j\), then
\(\var(x_1+\dots+x_m)=\var(x_1)+\dots+\var(x_m)\).
\item \(\expc(uu')=\var(u)+\mu_u\mu_u'\).
\end{enumerate}
\end{theorem}
\begin{proof}
\begin{board}[(a) By Theorem~\ref{h-thm-linear}(b): \(Au+b-\expc(Au+b)=A(u-\mu_u)\), \(Bv+d-\expc(Bv+d)=B(v-\mu_v)\);
transposes (\lec{1}, Theorem~8):]
\boardtop
\begin{align*}
\cov(Au+b,Bv+d)&=\expc\bigl(A(u-\mu_u)(v-\mu_v)'B'\bigr)=A\cov(u,v)B'
&&\text{(Theorem~\ref{h-thm-linear}(a))}
\end{align*}
\(u=v=z\): \(\cov(Az,Bz)=A\var(z)B'\).
\end{board}
(b) \(\expc(u_1+u_2)=\expc(u_1)+\expc(u_2)\), so
\((u_1+u_2-\expc(u_1+u_2))(v-\mu_v)'=(u_1-\expc(u_1))(v-\mu_v)'+(u_2-\expc(u_2))(v-\mu_v)'\)
(\lec{1}, Theorem~3(a)); take expectations entry by entry (\ref{f-lin}). Transpose for the second identity
(Definition~\ref{h-def-cov}).

(c) By (b) in both arguments,
\(\var(\sum_ix_i)=\cov(\sum_ix_i,\sum_jx_j)=\sum_i\sum_j\cov(x_i,x_j)\); only the
\(m\) terms with \(i=j\) remain.

(d) \((u-\mu_u)(u-\mu_u)'=uu'-u\mu_u'-\mu_uu'+\mu_u\mu_u'\) (\lec{1}, Theorems~3(a), 8); take expectations entry by entry, with
\(\expc(u\mu_u')=\mu_u\mu_u'=\expc(\mu_uu')\) (Theorem~\ref{h-thm-linear}(a)):
\(\var(u)=\expc(uu')-\mu_u\mu_u'\).
\end{proof}

\Needspace{18\baselineskip}
\begin{example}\label{h-exm-cov}
Let \(z\in\R^4\) be a random vector with
\[
\var(z)=\Sigma=\begin{pmatrix}4&0&4&0\\0&1&2&2\\4&2&9&6\\0&2&6&9\end{pmatrix}
\]
Find \(\cov(z,e)\) and \(\var(e)\) for \(e=z_3-z_1-2z_2\) and for \(f=z_4+2z_1+2z_2-2z_3\);
\(\var(z_1+z_2+z_3+z_4)\); the variances and the covariance of \(z_1+z_2\) and \(z_1-z_2\).
\end{example}
\textit{Solution.}
\begin{board}[\(e=a'z\), \(a=(-1,-2,1,0)'\); \(f=b'z\), \(b=(2,2,-2,1)'\); by Theorem~\ref{h-thm-cov}(a),
\(\cov(z,a'z)=\Sigma a\), \(\var(a'z)=a'\Sigma a\):]
\boardtop
\begin{align*}
\Sigma a&=(-4+0+4+0,\ 0-2+2+0,\ -4-4+9+0,\ 0-4+6+0)'=(0,0,1,2)'\\
\Sigma b&=(8+0-8+0,\ 0+2-4+2,\ 8+4-18+6,\ 0+4-12+9)'=(0,0,0,1)'
\end{align*}
So \(\cov(z,e)=(0,0,1,2)'\) and \(\var(e)=a'\Sigma a=0+0+1+0=1\): \(e\) is uncorrelated with
\(z_1\), \(z_2\); \(\cov(z,f)=(0,0,0,1)'\) and \(\var(f)=b'\Sigma b=1\): \(f\) is uncorrelated
with \(z_1\), \(z_2\), \(z_3\).
\end{board}
\begin{board}[\(z_1+z_2+z_3+z_4=\one'z\), \(z_1+z_2\), \(z_1-z_2\) (Theorem~\ref{h-thm-cov}(a)):]
\boardtop
\begin{align*}
\var(\one'z)&=\one'\Sigma\one=51\qquad(\text{the sum of all entries of }\Sigma)\\
(1,1,0,0)\Sigma&=(4,1,6,2),\qquad(1,-1,0,0)\Sigma=(4,-1,2,-2)
\end{align*}
So \(\var(z_1+z_2)=4+1=5\), \(\var(z_1-z_2)=4+1=5\), and
\(\cov(z_1+z_2,z_1-z_2)=(4,1,6,2)(1,-1,0,0)'=4-1=3\).
\end{board}

\Needspace{11\baselineskip}
\begin{theorem}\label{h-thm-quadratic}
Let \(z\in\R^n\) be a random vector, \(\expc(z)=\mu\), \(\var(z)=\Sigma\), and
\(A\in\R^{n\times n}\) constant. Then \(\expc(z'Az)=\tr(A\Sigma)+\mu'A\mu\); so
\(\expc(z'Az)=\tr(A\Sigma)\) if \(\mu=0\).
\end{theorem}
\begin{proof}
\begin{board}[\(x'Ax=\sum_i\sum_ja_{ij}x_ix_j\) for \(x=z\) and \(x=\mu\) (\lec{8}, Definition~12); \(\Sigma\) is symmetric
(Definition~\ref{h-def-random}); \(\tr(A\Sigma)=\sum_i\sum_ja_{ij}\sigma_{ji}\) (\lec{1}, Definition~9, Theorem~4(d)):]
\boardtop
\begin{align*}
\expc(z'Az)&=\textstyle\sum_i\sum_ja_{ij}\expc(z_iz_j)&&\text{(\ref{f-lin})}\\
\sigma_{ij}&=\expc\bigl((z_i-\mu_i)(z_j-\mu_j)\bigr)=\expc(z_iz_j)-\mu_i\mu_j-\mu_j\mu_i+\mu_i\mu_j
&&\text{(\ref{f-prob}, \ref{f-lin})}\\
\expc(z'Az)&=\textstyle\sum_i\sum_ja_{ij}\sigma_{ij}+\sum_i\sum_ja_{ij}\mu_i\mu_j=\tr(A\Sigma)+\mu'A\mu
&&(\sigma_{ji}=\sigma_{ij})
\end{align*}
\end{board}
\nopagebreak
\end{proof}

\begin{definition}\label{h-def-estimator}
Let observations be modeled as a random vector \(y\) whose distribution depends on unknown parameters. An
\textbf{estimator} of a parameter \(\theta\) is a random variable \(T=g(y)\), where the function \(g\) does not involve
unknown parameters; \(T\) is \textbf{unbiased} if \(\expc(T)=\theta\) for all values of the parameters.
\end{definition}

\Needspace{6\baselineskip}
\begin{example}\label{h-exm-quadratic}
Let \(y\in\R^n\), \(\expc(y)=\mu\), \(\var(y)=\sigma^2I_n\), and \(P\in\R^{n\times n}\) with
\(P'=P=P^2\) and \(\rank P=r\). Find \(\expc(y'Py)\), and the expectation of \(s^2=\frac1{n-1}\sum_i(y_i-\bar y)^2\).
\end{example}
\textit{Solution.} \(\tr P=\rank P=r\) (\lec{8}, Theorem~11: symmetric \(A^2=A\Rightarrow\rank A=\tr A\)), and by
Theorem~\ref{h-thm-quadratic} and \lec{1}, Theorem~10(a),
\(\expc(y'Py)=\tr(P\sigma^2I_n)+\mu'P\mu=\sigma^2r+\mu'P\mu\).
\begin{itemize}
\item \(P=C_n=I_n-\frac1n\one\one'\) (\(n\geqslant2\); \lec{1}, Definition~14): \(C_n'=C_n=C_n^2\), \(C_n\one=0\)
(\lec{1}, Proposition~15), \(\rank C_n=n-1\) (\lec{8}, Example~5); \(y'C_ny=\sum_i(y_i-\bar y)^2\) and
\(\mu'C_n\mu=\sum_i(\mu_i-\bar\mu)^2\) (\(\bar y=\frac1n\one'y\); \lec{1}, Example~16). So
\(\expc\bigl(\sum_i(y_i-\bar y)^2\bigr)=(n-1)\sigma^2+\sum_i(\mu_i-\bar\mu)^2\).
\item Common mean, \(\mu=\mu_0\one\): \(C_n\mu=0\), and \(\expc(s^2)=\sigma^2\) for all \(\mu_0\), \(\sigma^2\):
\(s^2\) is an unbiased estimator of \(\sigma^2\) (Definition~\ref{h-def-estimator}). The estimator
\(\frac1n\sum_i(y_i-\bar y)^2\) has the expectation \(\frac{n-1}n\sigma^2<\sigma^2\) (\(\sigma^2>0\)): not unbiased.
\item Drifting means: for \(n=4\) and \(\mu=(1,3,5,7)'\), \(\bar\mu=4\) and \(\sum_i(\mu_i-\bar\mu)^2=9+1+1+9=20\), so
\(\expc(s^2)=\sigma^2+\frac{20}3\): \(s^2\) overestimates \(\sigma^2\) on average.
\end{itemize}

\subsection*{2. Data matrices and the sample covariance matrix}

\begin{definition}\label{h-def-data}
A \textbf{data matrix} is \(X=(x_{ij})\in\R^{n\times p}\) (\(n\geqslant2\)): \(x_{ij}\) is the value of variable \(j\) on
unit \(i\); row \(i\) of \(X\) is \(x_i'=(x_{i1},\dots,x_{ip})\). With \(\one=\one_n\) and \(C_n=I_n-\frac1n\one\one'\) (\lec{1}, Definition~14):
the \textbf{mean vector} \(\bar x=\frac1nX'\one=(\bar x_1,\dots,\bar x_p)'\), \(\bar x_j=\frac1n\sum_ix_{ij}\); the
\textbf{centered data matrix} \(C_nX=X-\one\bar x'\) (\lec{1}, Example~17), with entries \(x_{ij}-\bar x_j\); the \textbf{sample covariance
matrix} \(S=(s_{jl})=\frac1{n-1}X'C_nX\in\R^{p\times p}\). For \(y,w\in\R^n\),
\(\frac1{n-1}\sum_i(y_i-\bar y)(w_i-\bar w)\) is their \textbf{sample covariance} (\(y\), \(w\) are
\textbf{uncorrelated} if it is \(0\)), and \(\frac1{n-1}\sum_i(y_i-\bar y)^2\) is the \textbf{sample variance} of \(y\).
For \(a\in\R^p\), \(Xa=(a'x_1,\dots,a'x_n)'\) holds the values on the \(n\) units of the linear combination of the
variables with coefficient vector \(a\).
\end{definition}

\Needspace{16\baselineskip}
\begin{proposition}\label{h-prop-data}
Let \(X\in\R^{n\times p}\) be a data matrix, with \(\bar x\) and \(S\) as in Definition~\ref{h-def-data}.
\begin{enumerate}
\item \(X'C_nX=(C_nX)'(C_nX)=X'X-n\bar x\bar x'\); \(s_{jj}\) is the sample variance of variable \(j\) and \(s_{jl}\)
the sample covariance of variables \(j\) and \(l\); \(S\) is symmetric and PSD.
\item For \(a,b\in\R^p\): \(Xa\) has mean \(a'\bar x\) and centered values \(C_n(Xa)=(C_nX)a\); the sample covariance
of \(Xa\) and \(Xb\) is \(a'Sb\), and the sample variance of \(Xa\) is \(a'Sa\).
\end{enumerate}
\end{proposition}
\begin{proof}
\begin{board}[(a) \(C_n'=C_n=C_n^2\) (\lec{1}, Proposition~15), \(X'\one=n\bar x\); column \(j\) of \(C_nX\) holds
\(x_{ij}-\bar x_j\):]
\boardtop
\begin{align*}
X'C_nX&=X'C_n'C_nX=(C_nX)'(C_nX)&&\text{(\lec{1}, Theorem~8)}\\
X'C_nX&=X'X-\tfrac1n(X'\one)(\one'X)=X'X-n\bar x\bar x'&&\text{(\lec{1}, Theorems~3(a), 8)}\\
\bigl((C_nX)'(C_nX)\bigr)_{jl}&=\textstyle\sum_i(x_{ij}-\bar x_j)(x_{il}-\bar x_l)&&\text{(\lec{1}, Theorem~4(d))}
\end{align*}
Divide by \(n-1\): \(s_{jl}\), \(s_{jj}\). \(S\) is symmetric (\lec{1}, Theorem~8) and PSD (\lec{8}, Theorem~15(b):
\(A'A\) is PSD).
\end{board}
\begin{board}[(b) \(\frac1n\one'(Xa)=\bar x'a=a'\bar x\), \(C_n(Xa)=(C_nX)a\): the deviations of \(Xa\) from its mean
(\lec{1}, Theorem~8, Example~16):]
\boardtop
\begin{align*}
\tfrac1{n-1}(C_nXa)'(C_nXb)&=a'\bigl(\tfrac1{n-1}(C_nX)'(C_nX)\bigr)b=a'Sb&&\text{(\lec{1}, Theorem~8; (a))}
\end{align*}
\(b=a\): the sample variance of \(Xa\) is \(a'Sa\).
\end{board}
\nopagebreak
\end{proof}

\Needspace{18\baselineskip}
\begin{example}\label{h-exm-data}
Five units (\(n=5\)) are measured on four variables (\(p=4\)), for instance five students on four tests:
\[
X=\begin{pmatrix}7&3&7&3\\9&9&5&7\\5&7&9&9\\3&5&1&1\\1&1&3&5\end{pmatrix}
\]
Find \(\bar x\), \(C_5X\) and \(S\).
\end{example}
\textit{Solution.}
\begin{board}[\(\one'X\), \(\bar x\); subtract \(\bar x'\) from every row (Definition~\ref{h-def-data}):]
\[
\one'X=(25,\ 25,\ 25,\ 25),\qquad\bar x=\tfrac15X'\one=\begin{pmatrix}5\\5\\5\\5\end{pmatrix},\qquad
C_5X=\begin{pmatrix}2&-2&2&-2\\4&4&0&2\\0&2&4&4\\-2&0&-4&-4\\-4&-4&-2&0\end{pmatrix}
\]
\end{board}
\begin{board}[\((C_5X)'(C_5X)\): inner products of the columns (\lec{1}, Theorem~4(d)); \(S\) (Proposition~\ref{h-prop-data}(a)):]
\[
(C_5X)'(C_5X)=\begin{pmatrix}40&28&20&12\\28&40&12&20\\20&12&40&28\\12&20&28&40\end{pmatrix},\qquad
S=\frac14(C_5X)'(C_5X)=\begin{pmatrix}10&7&5&3\\7&10&3&5\\5&3&10&7\\3&5&7&10\end{pmatrix}
\]
For instance \(40=4+16+0+4+16\) (column \(1\)) and \(28=-4+16+0+0+16\) (columns \(1\), \(2\)). Every variable has the
sample variance \(10\); the pairs \(\{1,2\}\), \(\{3,4\}\) have the sample covariance \(7\), the pairs \(\{1,3\}\),
\(\{2,4\}\) have \(5\), and \(\{1,4\}\), \(\{2,3\}\) have \(3\).
\end{board}

\Needspace{9\baselineskip}
\begin{theorem}\label{h-thm-sample}
Let \(x_1,\dots,x_n\in\R^p\) (\(n\geqslant2\)) be random vectors with \(\expc(x_i)=\mu\),
\(\var(x_i)=\Sigma\) and \(\cov(x_i,x_j)=0\) for \(i\ne j\); \(X\) the random data matrix with
rows \(x_1',\dots,x_n'\), \(\bar x=\frac1nX'\one\) and \(S=\frac1{n-1}X'C_nX\).
\begin{enumerate}
\item \(\expc(\bar x)=\mu\) and \(\var(\bar x)=\frac1n\Sigma\).
\item \(\expc(S)=\Sigma\).
\end{enumerate}
\end{theorem}
\begin{proof}
(a) \(X'\one=x_1+\dots+x_n\) (\lec{1}, Theorem~4(b)), so \(\expc(\bar x)=\frac1n(\mu+\dots+\mu)=\mu\)
(Theorem~\ref{h-thm-linear}(a)), and \(\var(\bar x)=\frac1{n^2}\var(x_1+\dots+x_n)
=\frac1{n^2}n\Sigma=\frac1n\Sigma\) (Theorem~\ref{h-thm-linear}(b) with \(A=\frac1nI_p\), Theorem~\ref{h-thm-cov}(c)).

(b) \(\expc(x_ix_i')=\Sigma+\mu\mu'\) and \(\expc(\bar x\bar x')=\frac1n\Sigma+\mu\mu'\)
(Theorem~\ref{h-thm-cov}(d); (a)). By Proposition~\ref{h-prop-data}(a) and \lec{1}, Theorem~4(d),
\(X'C_nX=X'X-n\bar x\bar x'\) with \(X'X=x_1x_1'+\dots+x_nx_n'\); so
\(\expc(X'C_nX)=n(\Sigma+\mu\mu')-n(\frac1n\Sigma+\mu\mu')=(n-1)\Sigma\), and \(\expc(S)=\Sigma\).
\end{proof}
With the divisor \(n\) in place of \(n-1\), the expectation would be \(\frac{n-1}n\Sigma\), not \(\Sigma\) (if
\(\Sigma\ne0\)); this is why \(S\) has the divisor \(n-1\).

\Needspace{6\baselineskip}
\begin{definition}\label{b9.17}
Let \(X\in\R^{n\times p}\) be a data matrix with every \(s_{jj}>0\); \(\sqrt{s_{jj}}\) is the \textbf{sample standard
deviation} of variable \(j\). With \(D_s=\diag(s_{11},\dots,s_{pp})\) and
\(D_s^{-1/2}=\diag\bigl(1/\sqrt{s_{11}},\dots,1/\sqrt{s_{pp}}\bigr)\), \(r_{jl}=s_{jl}/\sqrt{s_{jj}s_{ll}}\) is the
\textbf{sample correlation} of variables \(j\) and \(l\), and \(R=(r_{jl})=D_s^{-1/2}SD_s^{-1/2}\) (\lec{1},
Proposition~13(a)) is the \textbf{sample correlation matrix}.
\end{definition}

\Needspace{8\baselineskip}
\begin{proposition}\label{b9.18}
In Definition~\ref{b9.17}, \(r_{jj}=1\) and \(\lvert r_{jl}\rvert\leqslant1\).
\end{proposition}
\begin{proof}
\begin{board}[\(u_j=\col_j(C_nX)\): \((n-1)s_{jl}=u_j'u_l\) (Proposition~\ref{h-prop-data}(a)), \(\lVert u_j\rVert^2=u_j'u_j\):]
\[
r_{jj}=\frac{s_{jj}}{s_{jj}}=1,\qquad r_{jl}=\frac{u_j'u_l}{\lVert u_j\rVert\,\lVert u_l\rVert},\qquad\lvert
r_{jl}\rvert\leqslant1\qquad\text{(\lec{6}, Theorem~3: \(\lvert x'y\rvert\leqslant\lVert x\rVert\,\lVert y\rVert\))}
\]
\end{board}
\nopagebreak
\end{proof}
\begin{board}[Example~\ref{h-exm-data}: every \(s_{jj}=10\), so \(D_s^{-1/2}=\frac1{\sqrt{10}}I_4\) and \(R=\frac1{10}S\):]
\(r_{12}=s_{12}/\sqrt{s_{11}s_{22}}=7/\sqrt{10\cdot10}=\frac7{10}\), \(r_{13}=\frac5{10}=\frac12\),
\(r_{14}=\frac3{10}\).
\end{board}

\Needspace{19\baselineskip}
\subsection*{3. The Rayleigh quotient}

\begin{theorem}\label{h-thm-rayleigh}
Let \(A\in\R^{n\times n}\) be symmetric, and let \(q_1,\dots,q_n\) be orthonormal (\lec{6}, Definition~8) eigenvectors of \(A\),
\(Aq_i=\lambda_iq_i\), numbered so that \(\lambda_1\geqslant\dots\geqslant\lambda_n\) (\lec{8}, Theorem~9(a)); then
\(Q=[q_1\ \cdots\ q_n]\) is orthogonal, \(Q'AQ=\diag(\lambda_1,\dots,\lambda_n)\), and the \(\lambda_i\) are the
eigenvalues of \(A\) (\lec{8}, Theorem~9(b)). For \(x\ne0\), \(x'Ax/x'x\) is the \textbf{Rayleigh quotient} of \(x\).
\begin{enumerate}
\item \(\lambda_n\leqslant x'Ax/x'x\leqslant\lambda_1\) for every \(x\ne0\), with equality at \(x=q_n\) and \(x=q_1\):
\(\max_{x\ne0}x'Ax/x'x=\lambda_1\) and \(\min_{x\ne0}x'Ax/x'x=\lambda_n\).
\item For \(x\ne0\): \(x'Ax=\lambda_1x'x\iff Ax=\lambda_1x\), and \(x'Ax=\lambda_nx'x\iff Ax=\lambda_nx\).
\end{enumerate}
\end{theorem}
\begin{proof}
\begin{board}[(a) \(y=Q'x\): \(x=QQ'x=Qy\) (\(QQ'=I\), \lec{6}, Theorem~9(b)):]
\boardtop
\begin{align*}
x'x&=y'Q'Qy=y'y,\qquad x'Ax=\lambda_1y_1^2+\dots+\lambda_ny_n^2&&\text{(\lec{8}, Theorem~14)}\\
\lambda_nx'x&=\textstyle\sum_i\lambda_ny_i^2\leqslant\sum_i\lambda_iy_i^2\leqslant\sum_i\lambda_1y_i^2=\lambda_1x'x
&&(\lambda_n\leqslant\lambda_i\leqslant\lambda_1)\\
q_1'Aq_1&=\lambda_1q_1'q_1=\lambda_1,\qquad q_n'Aq_n=\lambda_nq_n'q_n=\lambda_n
\end{align*}
Divide by \(x'x>0\).
\end{board}
\begin{board}[(b) \(x\ne0\), \(y=Q'x\), \(x=Qy=\sum_iy_iq_i\) (\lec{1}, Theorem~4(b)):]
\boardtop
\begin{align*}
\lambda_1x'x-x'Ax&=\textstyle\sum_i(\lambda_1-\lambda_i)y_i^2\geqslant0\\
\lambda_1x'x-x'Ax=0&\iff y_i=0\text{ for }\lambda_i<\lambda_1\implies x=\textstyle\sum_{\lambda_i=\lambda_1}y_iq_i
\implies Ax=\lambda_1x\\
Ax=\lambda_1x&\implies x'Ax=\lambda_1x'x
\end{align*}
For \(\lambda_n\), apply this to \(-A\): \((-A)q_i=-\lambda_iq_i\), and in the order \(q_n,\dots,q_1\) the eigenvalues
\(-\lambda_n\geqslant\dots\geqslant-\lambda_1\) of \(-A\) decrease.
\end{board}
\nopagebreak
\end{proof}

\Needspace{17\baselineskip}
\begin{theorem}\label{h-thm-constrained}
Let \(A\), \(q_1,\dots,q_n\) and \(\lambda_1\geqslant\dots\geqslant\lambda_n\) be as in Theorem~\ref{h-thm-rayleigh},
and \(1\leqslant k\leqslant n\). Then
\[
\lambda_k=\max\Bigl\{\frac{x'Ax}{x'x}:x\ne0,\ x\perp q_1,\dots,q_{k-1}\Bigr\}
\]
attained at \(x=q_k\) (for \(k=1\) there is no condition); \(x\perp q\) means \(q'x=0\) (\lec{6}, Definition~1).
\end{theorem}
\begin{proof}
\begin{board}[\(Q=[R\ T]\), \(R=[q_1\ \cdots\ q_{k-1}]\), \(T=[q_k\ \cdots\ q_n]\); \(QQ'=I\) (\lec{6}, Theorem~9(b));
\(AT=[\lambda_kq_k\ \cdots\ \lambda_nq_n]=T\diag(\lambda_k,\dots,\lambda_n)\) (\lec{1}, Theorem~4(a), Proposition~13(a)):]
\boardtop
\begin{align*}
&R'T=0,\ T'T=I_{n-k+1},\ I=QQ'=RR'+TT'&&\text{(\lec{1}, Theorem~4(d))}\\
x\perp q_1,\dots,q_{k-1}&\iff R'x=0\iff x=Tz,\ z=T'x&&\text{(\lec{1}, Theorem~4(c); line 1)}\\
x'x&=z'T'Tz=z'z,\qquad x'Ax=z'(T'AT)z\\
T'AT&=T'T\diag(\lambda_k,\dots,\lambda_n)=\diag(\lambda_k,\dots,\lambda_n)
\end{align*}
By Theorem~\ref{h-thm-rayleigh}(a) for \(\diag(\lambda_k,\dots,\lambda_n)\) (with \(Q=I\)), the max of
\(z'(T'AT)z/z'z\) over \(z\ne0\) is \(\lambda_k\), at \(z=e_1\): \(x=Te_1=q_k\).
\end{board}
\nopagebreak
\end{proof}

\subsection*{4. Principal components}

\begin{definition}\label{h-def-pc}
Let \(X\in\R^{n\times p}\) be a data matrix with \(S\ne0\). A \textbf{first loading vector} is a unit vector
\(v_1\in\R^p\) that maximizes \(a'Sa\) over all unit vectors \(a\in\R^p\) (\(\lVert a\rVert=1\), \lec{6}, Definition~1). For \(m\geqslant2\), once \(v_1,\dots,v_{m-1}\)
are chosen, an \textbf{\(m\)th loading vector} is a unit vector \(v_m\) that maximizes \(a'Sa\) over the unit vectors
\(a\) with \(a'Sv_l=0\) (\(l=1,\dots,m-1\)) (the \textbf{\(m\)th problem}). The linear combination with coefficient
vector \(v_m\) is the \textbf{\(m\)th principal component}, and \(z_m=C_nXv_m\in\R^n\), with the entries
\(z_{im}=v_m'(x_i-\bar x)\), is the \textbf{\(m\)th score vector}.
\end{definition}
By Proposition~\ref{h-prop-data}(b), \(a'Sa\) is the sample variance of \(Xa\), and \(a'Sv_l=0\) says that \(Xa\) and
\(Xv_l\) are uncorrelated: the \(m\)th principal component has the largest sample variance among the combinations with
a unit coefficient vector that are uncorrelated with the first \(m-1\) principal components.

\Needspace{11\baselineskip}
\begin{theorem}\label{h-thm-pca}
Let \(X\in\R^{n\times p}\) be a data matrix with \(S\ne0\), and \(S=V\Lambda V'\) with \(V=[v_1\ \cdots\ v_p]\)
orthogonal, \(\Lambda=\diag(\lambda_1,\dots,\lambda_p)\), \(\lambda_1\geqslant\dots\geqslant\lambda_p\) (as in
Theorem~\ref{h-thm-rayleigh}). Let \(r=\rank(C_nX)\) and \(z_m=C_nXv_m\).
\begin{enumerate}
\item \(\lambda_p\geqslant0\), \(r\leqslant\min\{n-1,p\}\), and \(\lambda_m>0\iff m\leqslant r\).
\item If \(\lambda_l>0\), then for every \(a\in\R^p\): \(Xa\) and \(Xv_l\) are uncorrelated \(\iff\) \(a\perp v_l\).
\item Let \(m\leqslant r\). If \(v_1,\dots,v_{m-1}\) are the first \(m-1\) loading vectors, then \(v_m\) is an \(m\)th
loading vector, and the maximum in the \(m\)th problem is \(\lambda_m\).
\item \(z_m'z_l=(n-1)\lambda_m\) if \(l=m\), and \(0\) if \(l\ne m\): the score vectors are pairwise orthogonal and
uncorrelated, and \(z_m\) has the sample variance \(\lambda_m\).
\end{enumerate}
\end{theorem}
\begin{proof}
(a)
\begin{itemize}
\item \(S\) is PSD (Proposition~\ref{h-prop-data}(a)), so \(\lambda_p\geqslant0\) (\lec{8}, Theorem~14(a)).
\item \(r=\rank(C_nX)\leqslant\rank C_n=n-1\) (\lec{4}, Theorem~15(a); \lec{8}, Example~5), \(r\leqslant p\) (\lec{4},
Theorem~12(a)).
\item \(\rank S=\rank\bigl((C_nX)'(C_nX)\bigr)=r\) (\lec{4}, Theorems~15(d), 16(b)) and \(\rank S=\#\{m:\lambda_m\ne0\}\)
(\lec{8}, Theorem~9(b)). As \(\lambda_1\geqslant\dots\geqslant\lambda_p\geqslant0\), exactly \(\lambda_1,\dots,\lambda_r\) are
positive.
\end{itemize}
(b) By Proposition~\ref{h-prop-data}(b) and \(Sv_l=\lambda_lv_l\), the sample covariance of \(Xa\) and \(Xv_l\) is
\(a'Sv_l=\lambda_la'v_l\), which is \(0\iff a'v_l=0\), as \(\lambda_l>0\).
\begin{board}[(c) Unit \(a\), \(c=V'a\): \(a=Vc\), \(c_j=v_j'a\), \(\sum_jc_j^2=\lVert c\rVert^2=\lVert a\rVert^2=1\)
(\lec{6}, Theorem~9(a),(b)); by (b), (a), the \(m\)th problem asks \(c_1=\dots=c_{m-1}=0\):]
\boardtop
\begin{align*}
a'Sa&=c'V'SVc=c'\Lambda c=\textstyle\sum_j\lambda_jc_j^2\\
\lambda_m-a'Sa&=\textstyle\sum_{j\geqslant m}\lambda_mc_j^2-\sum_{j\geqslant m}\lambda_jc_j^2
=\sum_{j>m}(\lambda_m-\lambda_j)c_j^2\geqslant0\\
v_m'Sv_l&=\lambda_lv_m'v_l=0\ (l<m),\qquad v_m'Sv_m=\lambda_m
\end{align*}
So \(v_m\) solves the \(m\)th problem, with the maximum \(\lambda_m\).
\end{board}
\begin{board}[(d) Proposition~\ref{h-prop-data}(a); \(\frac1n\one'z_m=\frac1n\one'C_nXv_m=0\) (\lec{1}, Proposition~15):]
\boardtop
\begin{align*}
z_m'z_l&=v_m'(C_nX)'(C_nX)v_l=(n-1)v_m'Sv_l=(n-1)\lambda_lv_m'v_l
\end{align*}
The score vectors have mean \(0\), so \(\frac1{n-1}z_m'z_l\) is their sample covariance (Definition~\ref{h-def-data}).
\end{board}
\nopagebreak
\end{proof}
Part (c) is Theorem~\ref{h-thm-constrained} for \(S\).

\Needspace{18\baselineskip}
\begin{proposition}\label{h-prop-pve}
In Theorem~\ref{h-thm-pca}, the \textbf{total variance} is \(s_{11}+\dots+s_{pp}=\tr S=\lambda_1+\dots+\lambda_p\),
and the \textbf{proportion of variance explained} by the \(m\)th principal component is
\[
\mathrm{PVE}_m=\frac{\lVert z_m\rVert^2}{\lVert C_nX\rVert_F^2}=\frac{\lambda_m}{\lambda_1+\dots+\lambda_p}\qquad
(m=1,\dots,p)
\]
where \(\lVert C_nX\rVert_F^2=\tr\bigl((C_nX)'(C_nX)\bigr)\) is the sum of the squared entries of \(C_nX\)
(\lec{1}, Theorem~10(c)). \(\mathrm{PVE}_m>0\) for \(m\leqslant r\), \(\mathrm{PVE}_m=0\) for \(m>r\),
\(\mathrm{PVE}_1\geqslant\mathrm{PVE}_2\geqslant\cdots\), and \(\mathrm{PVE}_1+\dots+\mathrm{PVE}_M\) increases with
\(M\leqslant r\) and equals \(1\) for \(M\geqslant r\).
\end{proposition}
\begin{proof}
\begin{board}[\(V'V=I\) (\lec{6}, Definition~8); Proposition~\ref{h-prop-data}(a), Theorem~\ref{h-thm-pca}(d):]
\boardtop
\begin{align*}
\tr S&=\tr(V\Lambda V')=\tr(\Lambda V'V)=\tr\Lambda=\lambda_1+\dots+\lambda_p&&\text{(\lec{1}, Theorem~10(b))}\\
\lVert C_nX\rVert_F^2&=\tr\bigl((C_nX)'(C_nX)\bigr)=(n-1)\tr S,\qquad\lVert z_m\rVert^2=(n-1)\lambda_m
\end{align*}
\(\tr S\geqslant\lambda_1>0\) (Theorem~\ref{h-thm-pca}(a); \(\Lambda\ne0\) as \(S=V\Lambda V'\ne0\)); \(\lambda_m>0\) exactly for
\(m\leqslant r\), and \(\lambda_m\) decreases: \(\mathrm{PVE}_1+\dots+\mathrm{PVE}_r=(\lambda_1+\dots+\lambda_r)/\tr S=1\).
\end{board}
\nopagebreak
\end{proof}

\Needspace{5\baselineskip}
\begin{example}\label{h-exm-pca}
Find the principal components, the score vectors and the proportions of variance explained for the data of
Example~\ref{h-exm-data}.
\end{example}
\par\Needspace{12\baselineskip}
\textit{Solution.} \(h_1=\one_4\), \(h_2=(1,1,-1,-1)'\), \(h_3=(1,-1,1,-1)'\), \(h_4=(1,-1,-1,1)'\).
\begin{board}[\(Sh_m\) row by row (\(S\) of Example~\ref{h-exm-data}); the \(h_m\) are orthogonal, \(\lVert h_m\rVert=2\):]
\[
Sh_1=\begin{pmatrix}25\\25\\25\\25\end{pmatrix}=25h_1,\quad
Sh_2=\begin{pmatrix}9\\9\\-9\\-9\end{pmatrix}=9h_2,\quad
Sh_3=\begin{pmatrix}5\\-5\\5\\-5\end{pmatrix}=5h_3,\quad
Sh_4=\begin{pmatrix}1\\-1\\-1\\1\end{pmatrix}=h_4
\]
For instance \(10+7+5+3=25\), \(10+7-5-3=9\), \(10-7+5-3=5\), \(10-7-5+3=1\) in row \(1\). So \(v_m=\frac12h_m\) are
orthonormal eigenvectors, the eigenvalues are \(\lambda_1=25\), \(\lambda_2=9\), \(\lambda_3=5\), \(\lambda_4=1\), and
\(S=V\diag(25,9,5,1)V'\) with \(V=[v_1\ \cdots\ v_4]\) (\lec{8}, Theorem~9(b)); \(r=4\)
(Theorem~\ref{h-thm-pca}(a)).
\end{board}
By Theorem~\ref{h-thm-pca}(c), the first principal component, with coefficient vector \(\frac12(1,1,1,1)'\) and sample
variance \(25\), is half the sum of the four variables: an overall level. The second, \(\frac12(1,1,-1,-1)'\), with
sample variance \(9\), contrasts variables \(1\), \(2\) with variables \(3\), \(4\); the third, \(\frac12(1,-1,1,-1)'\),
with sample variance \(5\), contrasts variables \(1\), \(3\) with \(2\), \(4\); the fourth, \(\frac12(1,-1,-1,1)'\), has
sample variance \(1\).
\begin{board}[Scores \(z_m=C_5Xv_m=\frac12C_5Xh_m\) (\(C_5X\) of Example~\ref{h-exm-data}); check (Theorem~\ref{h-thm-pca}(d)):]
\[
C_5Xh_1=\begin{pmatrix}0\\10\\10\\-10\\-10\end{pmatrix},\quad C_5Xh_2=\begin{pmatrix}0\\6\\-6\\6\\-6\end{pmatrix},\quad
C_5Xh_3=\begin{pmatrix}8\\-2\\-2\\-2\\-2\end{pmatrix},\quad C_5Xh_4=\begin{pmatrix}0\\2\\-2\\-2\\2\end{pmatrix}
\]
\[
z_1=\begin{pmatrix}0\\5\\5\\-5\\-5\end{pmatrix},\quad z_2=\begin{pmatrix}0\\3\\-3\\3\\-3\end{pmatrix},\quad
z_3=\begin{pmatrix}4\\-1\\-1\\-1\\-1\end{pmatrix},\quad z_4=\begin{pmatrix}0\\1\\-1\\-1\\1\end{pmatrix}
\]
\(\lVert z_1\rVert^2=100=4\cdot25\), \(\lVert z_2\rVert^2=36=4\cdot9\), \(\lVert z_3\rVert^2=20=4\cdot5\),
\(\lVert z_4\rVert^2=4=4\cdot1\), and \(z_m'z_l=0\) for \(m<l\); for instance \(z_1'z_2=0+15-15-15+15=0\) and
\(z_1'z_3=0-5-5+5+5=0\).
\end{board}
\begin{board}[Proportions (Proposition~\ref{h-prop-pve}):]
\(\tr S=10+10+10+10=40=25+9+5+1\), so \(\mathrm{PVE}_1=\frac{25}{40}=\frac58\), \(\mathrm{PVE}_2=\frac9{40}\),
\(\mathrm{PVE}_3=\frac5{40}=\frac18\), \(\mathrm{PVE}_4=\frac1{40}\): the first component explains \(\frac58\) of the
total variance, the first two together \(\frac{34}{40}=\frac{17}{20}\).
\end{board}

\begin{remark}\label{h-rem-number}
Usually only the first \(M\) principal components are kept, with \(M\) much smaller than \(p\); their score vectors
\(z_1,\dots,z_M\) summarize the data, for instance in a plot of the units. By Proposition~\ref{h-prop-pve}, they explain
together the proportion \(\mathrm{PVE}_1+\dots+\mathrm{PVE}_M\) of the total variance, and
\(\mathrm{PVE}_1\geqslant\mathrm{PVE}_2\geqslant\cdots\). A common aid is the \textbf{scree plot}, the points
\((m,\mathrm{PVE}_m)\), \(m=1,\dots,p\), joined by segments: one takes the smallest \(M\) that explains a sizable part of
the total variance, typically the \(M\) after which the \(\mathrm{PVE}_m\) drop off and stay small (an \textbf{elbow} of
the plot after the \(M\)th component). There is no objective rule; the choice depends on the data and on the purpose.
\end{remark}

\subsection*{5. Principal components of a random vector}

\begin{definition}\label{h-def-pop-pc}
Let \(x=(x_1,\dots,x_p)'\) be a random vector with \(\Sigma=\var(x)\ne0\). A \textbf{first principal
component} of \(x\) is \(a_1'x\), where the unit vector \(a_1\in\R^p\) maximizes \(\var(a'x)\) over all unit
vectors \(a\in\R^p\). For \(m\geqslant2\), once \(a_1,\dots,a_{m-1}\) are chosen, an \textbf{\(m\)th principal
component} is \(a_m'x\), where the unit vector \(a_m\) maximizes \(\var(a'x)\) over the unit vectors
\(a\) for which \(a'x\) is uncorrelated with each of \(a_1'x,\dots,a_{m-1}'x\).
\end{definition}

\Needspace{14\baselineskip}
\begin{theorem}\label{h-thm-pop-pca}
Let \(x\) and \(\Sigma\) be as in Definition~\ref{h-def-pop-pc}, and \(\Sigma=V\Lambda V'\) with
\(V=[v_1\ \cdots\ v_p]\) orthogonal, \(\Lambda=\diag(\lambda_1,\dots,\lambda_p)\),
\(\lambda_1\geqslant\dots\geqslant\lambda_p\) (as in Theorem~\ref{h-thm-rayleigh}).
\begin{enumerate}
\item \(\var(a'x)=a'\Sigma a\) and \(\cov(a'x,b'x)=a'\Sigma b\) for \(a,b\in\R^p\); and
\(\lambda_p\geqslant0\).
\item \(v_1'x\) is a first principal component, with variance \(\lambda_1\). Let \(m\geqslant2\) and \(\lambda_{m-1}>0\).
If \(v_1'x,\dots,v_{m-1}'x\) are the first \(m-1\) principal components, then \(v_m'x\) is an \(m\)th principal
component, with variance \(\lambda_m\).
\item \(v_1'x,\dots,v_p'x\) are pairwise uncorrelated, and the \textbf{total variance}
\(\var(x_1)+\dots+\var(x_p)=\tr\Sigma=\lambda_1+\dots+\lambda_p\) is positive. So the first
\(M\) principal components account for the proportion \((\lambda_1+\dots+\lambda_M)/\tr\Sigma\) of the total variance.
\end{enumerate}
\end{theorem}
\begin{proof}
(a) Let \(T\in\R^{2\times p}\) have the rows \(a'\), \(b'\), so \(Tx=(a'x,b'x)'\) and
\(\var(Tx)=T\Sigma T'\) (Theorem~\ref{h-thm-linear}(b)) has the entries \(a'\Sigma a\), \(a'\Sigma b\);
\(b'\Sigma a\), \(b'\Sigma b\) (\lec{1}, Theorem~4(c),(d)); by Definition~\ref{h-def-random}, its diagonal entries are
\(\var(a'x)\), \(\var(b'x)\) and its \((1,2)\) entry is \(\cov(a'x,b'x)\).
\(\Sigma\) is PSD (Theorem~\ref{h-thm-linear}(c)), so \(\lambda_p\geqslant0\) (\lec{8}, Theorem~14(a)).

\Needspace{4\baselineskip}
(b)
\begin{itemize}
\item By (a), the first problem maximizes \(a'\Sigma a\) over the unit \(a\); for \(a\ne0\), \(a'\Sigma a/a'a\) is the
value at the unit vector \(a/\lVert a\rVert\). By Theorem~\ref{h-thm-rayleigh}(a) the maximum is \(\lambda_1\), at
\(a=v_1\).
\item Let \(m\geqslant2\) and \(\lambda_{m-1}>0\); then \(\lambda_l>0\) for \(l<m\). By (a) and
\(\Sigma v_l=\lambda_lv_l\), \(\cov(a'x,v_l'x)=a'\Sigma v_l=\lambda_la'v_l\), which is \(0\iff a\perp v_l\).
So the \(m\)th problem maximizes \(a'\Sigma a\) over the unit \(a\perp v_1,\dots,v_{m-1}\).
\item By Theorem~\ref{h-thm-constrained}, the maximum is \(\lambda_m\), at \(a=v_m\), and
\(\var(v_m'x)=v_m'\Sigma v_m=\lambda_m\).
\end{itemize}
(c) By (a), \(\cov(v_m'x,v_l'x)=v_m'\Sigma v_l=\lambda_lv_m'v_l=0\) for \(l\ne m\).
\(\var(x_j)=\sigma_{jj}\) (Definition~\ref{h-def-random}), so the total variance is
\(\tr\Sigma=\tr(V\Lambda V')=\tr(\Lambda V'V)=\tr\Lambda\) (\lec{1}, Theorem~10(b)), which is positive: every \(\lambda_m\geqslant0\)
by (a), and \(\Lambda\ne0\), as \(\Sigma=V\Lambda V'\ne0\).
\end{proof}
Parts (b) and (c) are Theorem~\ref{h-thm-pca}(c),(d) and Proposition~\ref{h-prop-pve} with \(\Sigma\) in place of \(S\).
In practice \(\Sigma\) is not known: \(n\) observations of \(x\) form the rows of a data matrix, and the principal
components of the data (Definition~\ref{h-def-pc}), computed from \(S\), take the place of those of \(x\).

\Needspace{14\baselineskip}
\begin{example}\label{h-exm-equicorrelation}
Let \(x=(x_1,\dots,x_p)'\) (\(p\geqslant2\)) have \(\var(x_i)=\sigma^2>0\) and
\(\cov(x_i,x_j)=\rho\sigma^2\) for \(i\ne j\), where \(0<\rho<1\):
\(\Sigma=\sigma^2\bigl((1-\rho)I_p+\rho J_p\bigr)\) (\(J_p=\one_p\one_p'\)). Find the principal components of \(x\).
\end{example}
\textit{Solution.}
\begin{board}[\(I_p=\bar J_p+C_p\), \(J_p=p\bar J_p\); for \(a\in\R^p\), \(\bar J_pa=\bar a\one\), \(C_pa=a-\bar a\one\)
(\lec{1}, Definition~14, Example~16):]
\boardtop
\begin{align*}
\Sigma&=\lambda_1\bar J_p+\lambda_2C_p,\qquad\lambda_1=\sigma^2\bigl(1+(p-1)\rho\bigr),\qquad\lambda_2=\sigma^2(1-\rho),\qquad\lambda_1-\lambda_2=p\rho\sigma^2>0\\
\Sigma\one&=\lambda_1\one,\qquad\Sigma a=\lambda_2a\ \text{ if }\one'a=0\\
\Sigma-\lambda_1I_p&=(\lambda_2-\lambda_1)C_p,\qquad\text{so }\Sigma a=\lambda_1a\iff C_pa=0\iff a=\bar a\one
\end{align*}
\(v_1=\frac1{\sqrt p}\one\); \([v_1\ \cdots\ v_p]\) is orthogonal for some \(v_2,\dots,v_p\) (\lec{6},
Corollary~11(c)); then \(\one'v_j=\sqrt p\,v_1'v_j=0\), \(\Sigma v_j=\lambda_2v_j\) (\(j\geqslant2\)): the eigenvalues
are \(\lambda_1\) and \(\lambda_2\) (\(p-1\) times) (\lec{8}, Theorem~9(b)); \(\lambda_1>\lambda_2>0\).
\end{board}
\begin{board}[The principal components (Theorem~\ref{h-thm-pop-pca}):]
\begin{itemize}
\item First: \(v_1'x=\frac1{\sqrt p}(x_1+\dots+x_p)\), with variance \(\lambda_1\). Unique up to sign: a unit \(a\) with
\(a'\Sigma a=\lambda_1\) has \(\Sigma a=\lambda_1a\) (Theorem~\ref{h-thm-rayleigh}(b)), so \(a=\bar a\one=\pm v_1\).
\item Second: \(\cov(a'x,v_1'x)=\lambda_1a'v_1\), so the constraint is \(\one'a=0\); every such unit \(a\)
has \(\Sigma a=\lambda_2a\) and \(\var(a'x)=\lambda_2\). Every \(a'x\) with \(\one'a=0\), \(\lVert a\rVert=1\)
is a second principal component: for \(p\geqslant3\) not unique, even up to sign (for \(p=2\) it is
\(\pm(x_1-x_2)/\sqrt2\)).
\item The first accounts for \(\lambda_1/\tr\Sigma=\bigl(1+(p-1)\rho\bigr)/p\) of \(\tr\Sigma=p\sigma^2\); for
\(p=4\), \(\rho=\frac12\): \(\frac58\).
\end{itemize}
\end{board}

\bigskip
\textbf{Practice.} The Exercises and Problems of the textbook, Chapter~9.

\end{document}
